

%%%%%%%%%%%%%%%
\begin{frame}[fragile]{Jointures}

Jointure =  \alert{produit cartésien} composé
avec une \alert{sélection}.

\vfill

\begin{minted}[baselinestretch=.9,
               fontsize=\small,
               framesep=2mm]{sql}
select expressions
from  table1 cross join table2
[where condition_jointure]
\end{minted}

\vfill

\alert{Interprétation}\,: \textit{comme avant}, l'espace de recherche (\sqlkw{from})
est maintenant la table résultat de $table_1 \times table_2$

\vfill
\begin{remblock}{autres syntaxes}
SQL propose plusieurs syntaxes équivalentes pour
la jointure. Nous donnons les autres options plus tard.
\end{remblock}

\end{frame}


%%%%%%%%%%%%%%%
\begin{frame}[fragile]{Cas extrême\,: le produit cartésien}

Sans clause \sqlkw{where}, on effectue un \alert{produit cartésien}.

\vfill

\begin{minted}[baselinestretch=.9,
               fontsize=\small,
               framesep=2mm]{sql}
select * from Logement cross join Activité
\end{minted}

\vfill
\begin{footnotesize}
\begin{tabular}{c|c|c|c|c|c}
\textbf{code} & \textbf{nom} & \textbf{capacité} & \textbf{type} & \textbf{codeLog.} & \textbf{codeActivité}\\ \hline 
ca & Causses & 45 & Auberge & ca & Randonnée\\ \hline 
ge & Génépi & 134 & Hôtel & ca & Randonnée\\ \hline 
pi & U Pinzutu & 10 & Gîte & ca & Randonnée\\ \hline 
ta & Tabriz & 34 & Hôtel & ca & Randonnée\\ \hline 
ca & Causses & 45 & Auberge & ge & Piscine\\ \hline 
ge & Génépi & 134 & Hôtel & ge & Piscine\\ \hline 
pi & U Pinzutu & 10 & Gîte & ge & Piscine\\ \hline 
... & ... & ... & ... & ... & ...\\ \hline 
\end{tabular}
\end{footnotesize}

\vfill

Quels nuplets nous intéressent particulièrement\,?

\end{frame}


%%%%%%%%%%%%%%%
\begin{frame}[fragile]{En ajoutant la sélection}

Le plus souvent, la jointure est ``naturelle''\,: la clé primaire
d'une table doit être égale à la clé étrangère d'une autre.
\vfill

\begin{minted}[baselinestretch=.9,
               fontsize=\small,
               framesep=2mm]{sql}
select * from Logement cross join Activité
where code = codeLogement
\end{minted}

\vfill
\begin{footnotesize}
\begin{tabular}{c|c|c|c|c|c}
\textbf{code} & \textbf{nom} & \textbf{capacité} & \textbf{type} & \textbf{codeLog.} & \textbf{codeActivité}\\ \hline 
ca & Causses & 45 & Auberge & ca & Randonnée\\ \hline 
ge & Génépi & 134 & Hôtel & ge & Piscine\\ \hline 
ge & Génépi & 134 & Hôtel & ge & Ski\\ \hline 
pi & U Pinzutu & 10 & Gîte & pi & Plongée\\ \hline 
pi & U Pinzutu & 10 & Gîte & pi & Voile\\ \hline 
\end{tabular}
\end{footnotesize}

\end{frame}

%%%%%%%%%%%%%%%
\begin{frame}[fragile]{Plus de deux tables\,?}

Aucun problème\,: on les énumère dans le \sqlkw{from}.

\vfill
Exemple\,: les voyageurs et les logements qu'ils ont visités.

\vfill

\begin{minted}[baselinestretch=.9,
               fontsize=\small,
               framesep=2mm]{sql}
select V.nom as nomVoyageur, L.nom as nomLogement
from Logement as L, Séjour as S, Voyageur as V
where code = S.codeLogement
and   S.idVoyageur = V.idVoyageur
\end{minted}

\vfill

\alert{Remarque 1}\,: dans la clause \sqlkw{from}, on utilise le plus souvent 
la virgule à la place de \sqlkw{cross join}. 

\vfill

\alert{Remarque 2}\,: le renommage \sqlkw{as} sert, notamment,
à  gérer les ambiguités soulevées par les noms d'attributs.

\end{frame}



%%%%%%%%%%%%%%%
\begin{frame}[fragile]{D'autres syntaxes\,?}

SQL connaît la jointure naturelle.

\vfill

\begin{minted}[baselinestretch=.9,
               fontsize=\small,
               framesep=2mm]{sql}
select *
from Séjour natural join Voyageur 
\end{minted}

\vfill
\begin{footnotesize}
\begin{tabular}{c|c|c|c|c|c|c|c}
\textbf{idV.} & \textbf{idS.} & \textbf{codeL.} & \textbf{début} & \textbf{fin} & \textbf{nom} & \textbf{prénom} & \textbf{ville} \\ \hline 
10 & 1 & pi & 20 & 20 & Fogg & Phileas & Ajaccio \\ \hline 
10 & 6 & pi & 10 & 12 & Fogg & Phileas & Ajaccio \\ \hline 
20 & 2 & ta & 21 & 22 & Bouvier & Nicolas & Aurillac \\ \hline 
20 & 4 & pi & 19 & 23 & Bouvier & Nicolas & Aurillac \\ \hline 
... & ... & ... & ... & ... & ... & ... & ... \\ \hline 
\end{tabular}
\end{footnotesize}

\vfill

NB\,: l'attribut \texttt{idVoyageur} apparaît une seule fois
\end{frame}


%%%%%%%%%%%%%%%
\begin{frame}[fragile]{D'autres syntaxes\,?}

SQL connaît aussi la jointure ``non naturelle'' (les attributs
de jointure n'ont pas le même nom).

\vfill

\begin{minted}[baselinestretch=.9,
               fontsize=\small,
               framesep=2mm]{sql}
select *
from Logement join Activité on (code=codeLogement) 
\end{minted}

\vfill

\alert{Intérêt}\,? Aucun. Connaître une syntaxe suffit.  Celle
que nous avons présenté au début est la plus lisible et la plus neutre.

\vfill

\begin{remblock}{La syntaxe ne compte pas\,!}
Il existe 100 façons d'écrire une même requête\,: aucune n'est meilleure
qu'une autre\,: le système décide seul
de l'évaluation\,!
\end{remblock}
\end{frame}


\begin{frame}{À retenir}

La jointure est obtenue en effectuant un produit cartésien dans le \sqlkw{from}.

\vfill

\begin{redbullet}
 \item Tout se passe comme si on interrogeait une seule table, celle
     définie par $table_1 \times table_2$
  \item Le critère de jointure est exprimé dans le \sqlkw{where},
       avec les critères (éventuels) de sélection
  \item Les synonymes (attributs ou tables) entrainent des ambiguités, 
      et nécessitent des renommages
\end{redbullet}

\vfill

Beaucoup de syntaxes alternatives (dues aux deux paradigmes fondant SQL). En connaître 
une suffit.

\end{frame}

